\documentclass[11pt]{article}
% Tables, as the UltraCanvas LaTeX document reader maps them: a booktabs
% table, spans in both directions (\multicolumn, \multirow), a tabularx
% whose X column takes the leftover width, and a longtable with a repeated
% head. Every cell is rich text, so formulas inside cells are typeset too.
\usepackage{amsmath}
\usepackage{booktabs}
\usepackage{tabularx}
\usepackage{longtable}
\usepackage{multirow}
\title{Tables of the Solar System}
\author{The UltraCanvas Team}
\date{September 2026}
\begin{document}
\maketitle

\begin{abstract}
Four tables, each written the way a paper writes it. The reader turns every
one of them into the same rich-document table the ODT and DOCX readers
produce, so they render, export and re-import identically.
\end{abstract}

\section{Rules instead of lines}
The \texttt{booktabs} rules --- \texttt{\textbackslash toprule},
\texttt{\textbackslash midrule}, \texttt{\textbackslash bottomrule} --- mark
the header row; the first row of every imported table is the header.

\begin{table}[h]
  \centering
  \caption{The four inner planets}\label{tab:inner}
  \begin{tabular}{lrrr}
    \toprule
    Planet & Radius (km) & Mass ($10^{24}$ kg) & Moons \\
    \midrule
    Mercury & 2\,440 & 0.330 & 0 \\
    Venus   & 6\,052 & 4.87  & 0 \\
    Earth   & 6\,371 & 5.97  & 1 \\
    Mars    & 3\,390 & 0.642 & 2 \\
    \bottomrule
  \end{tabular}
\end{table}

Table~\ref{tab:inner} needs no vertical rules. The classic form with
\texttt{|} column separators and \texttt{\textbackslash hline} is read just
as happily, as Table~\ref{tab:spans} shows.

\section{Spans in both directions}
\texttt{\textbackslash multicolumn} sets a cell's \emph{columnSpan} and
\texttt{\textbackslash multirow} its \emph{rowSpan}; a reader that ignored
either would shift every later cell one column to the left.

\begin{table}[h]
  \centering
  \caption{Orbits, grouped}\label{tab:spans}
  \begin{tabular}{|l|c|c|c|}
    \hline
    \multirow{2}{*}{Planet} & \multicolumn{2}{c|}{Orbit} & Day \\ \cline{2-3}
     & Period (a) & $e$ & (h) \\ \hline
    Jupiter & 11.86 & 0.049 & 9.9 \\ \hline
    Saturn  & 29.45 & 0.052 & 10.7 \\ \hline
    Uranus  & 84.02 & 0.047 & 17.2 \\ \hline
  \end{tabular}
\end{table}

\section{A column that takes the rest of the line}
In a \texttt{tabularx} the \texttt{X} column absorbs the width the fixed
columns leave over, which is how a notes column is written.

\begin{table}[h]
  \centering
  \caption{Named after}\label{tab:names}
  \begin{tabularx}{\textwidth}{lX}
    \toprule
    Planet & Note \\
    \midrule
    Mercury & The messenger of the Roman gods, for the speed with which the
              planet moves against the fixed stars. \\
    Neptune & Predicted from the perturbations of Uranus before it was ever
              observed, and found within a degree of the prediction. \\
    \bottomrule
  \end{tabularx}
\end{table}

\section{A table longer than a page}
\texttt{longtable} spreads over page breaks and repeats its head; the
importer keeps the head row and drops the pagination machinery, since a
rich document has no pages until it is printed. A long table is not a
float, so its caption is written with \texttt{\textbackslash captionof},
which names the counter to use.

\captionof{table}{Escape velocities}\label{tab:escape}
\begin{longtable}{lrr}
  \toprule
  Body & $v_{\text{esc}}$ (km/s) & $g$ (m/s$^2$) \\
  \midrule
  \endhead
  Sun     & 617.5 & 274.0 \\
  Jupiter & 59.5  & 24.8 \\
  Earth   & 11.2  & 9.81 \\
  Mars    & 5.0   & 3.72 \\
  Moon    & 2.4   & 1.62 \\
  Pluto   & 1.2   & 0.62 \\
  \bottomrule
\end{longtable}

The escape velocity in Table~\ref{tab:escape} is
\begin{equation}\label{eq:escape}
  v_{\text{esc}} = \sqrt{\frac{2GM}{r}},
\end{equation}
and \eqref{eq:escape} is why the value falls as a body's radius grows at
constant density.
\end{document}
